\documentclass[11pt]{article}

\usepackage{amsmath,amsthm,amssymb}
\usepackage{mathtools}
\usepackage{booktabs}
\usepackage[margin=1in]{geometry}

\newtheorem{theorem}{Theorem}
\newtheorem{lemma}[theorem]{Lemma}
\newtheorem{proposition}[theorem]{Proposition}
\newtheorem{corollary}[theorem]{Corollary}
\newtheorem{definition}[theorem]{Definition}

\theoremstyle{remark}
\newtheorem{remark}[theorem]{Remark}

\newcommand{\CC}{\mathbb{C}}
\newcommand{\ket}[1]{|#1\rangle}
\newcommand{\pisort}{\pi_{\mathrm{sort}}}

\title{The Operator Independence Theorem:\\
  $\pisort$ is not a function of $R(u)$}
\author{Clio}
\date{April 10, 2026}

\begin{document}
\maketitle

\begin{abstract}
We prove that the 0-Hecke sorting operator $\pisort$ and the five-vertex
Demazure $R$-matrix $R(u)$ are algebraically independent: no specialization,
limit, or conjugation relates them. The proof rests on a minimal polynomial
argument---$\pisort$ satisfies a degree-2 relation while $R(u)$ has minimal
polynomial of degree~4 for all $u \neq 0$. As a corollary, the three-level
integrability hierarchy (solvable $\supset$ quasi-solvable $\supset$ combinatorial)
consists of genuinely distinct algebraic levels, not degenerations of one another.
\end{abstract}

%% -----------------------------------------------------------------------
\section{Definitions}
%% -----------------------------------------------------------------------

\begin{definition}[Ice-rule $R$-matrix]
We work over $\CC$. Let $V = \CC^2$ with basis $\{\ket{0}, \ket{1}\}$,
and identify $V \otimes V \cong \CC^4$ with ordered basis
$\{\ket{00}, \ket{01}, \ket{10}, \ket{11}\}$.
A \emph{six-vertex} or \emph{ice-rule} $R$-matrix is a linear map
$R\colon V \otimes V \to V \otimes V$ preserving the total spin $i+j$ of
each basis vector $\ket{ij}$. In the standard basis it has the form
\[
  R = \begin{pmatrix}
    a_1 & 0   & 0   & 0   \\
    0   & b_1 & c_1 & 0   \\
    0   & c_2 & b_2 & 0   \\
    0   & 0   & 0   & a_2
  \end{pmatrix}.
\]
\end{definition}

\begin{definition}[Demazure five-vertex $R$-matrix]\label{def:Ru}
Set $b_2 = 0$ (five-vertex condition) and
\[
  a_1(u) = 1+u,\quad a_2 = 1,\quad b_1(u) = u,\quad
  b_2 = 0,\quad c_1 = 1,\quad c_2 = 1,
\]
giving the one-parameter family
\begin{equation}\label{eq:Ru}
  R(u) = \begin{pmatrix}
    1+u & 0 & 0 & 0 \\
    0   & u & 1 & 0 \\
    0   & 1 & 0 & 0 \\
    0   & 0 & 0 & 1
  \end{pmatrix}.
\end{equation}
\end{definition}

\begin{definition}[0-Hecke sorting operator]\label{def:pisort}
The \emph{sorting operator} $\pisort\colon V \otimes V \to V \otimes V$
is defined by
\[
  \pisort \ket{ij} =
  \begin{cases}
    \ket{01} & \text{if } (i,j) = (1,0), \\
    \ket{ij} & \text{otherwise},
  \end{cases}
\]
with matrix form
\begin{equation}\label{eq:pisort}
  \pisort = \begin{pmatrix}
    1 & 0 & 0 & 0 \\
    0 & 1 & 1 & 0 \\
    0 & 0 & 0 & 0 \\
    0 & 0 & 0 & 1
  \end{pmatrix}.
\end{equation}
\end{definition}

\begin{definition}[Braid relation]
A matrix $R \in \mathrm{End}(V \otimes V)$ satisfies the
\emph{constant braid relation} if
\begin{equation}\label{eq:braid}
  R_{12}\, R_{23}\, R_{12} = R_{23}\, R_{12}\, R_{23}
\end{equation}
in $\mathrm{End}(V^{\otimes 3})$, where $R_{12} = R \otimes I$ and
$R_{23} = I \otimes R$.
\end{definition}

%% -----------------------------------------------------------------------
\section{The Operator Independence Theorem}
%% -----------------------------------------------------------------------

\begin{theorem}[Operator Independence]\label{thm:main}
Let $R(u)$ and $\pisort$ be as in \eqref{eq:Ru} and \eqref{eq:pisort}. Then:
\begin{enumerate}
  \item \textbf{$\pisort$ is a 0-Hecke generator:}
    $\pisort^2 = \pisort$ and $\pisort$ satisfies~\eqref{eq:braid}.
  \item \textbf{$R(u)$ is not a Hecke generator:}
    For all $u \neq 0$, $R(u)$ has at least three distinct eigenvalues
    and satisfies no quadratic relation $(R - \alpha I)(R - \beta I) = 0$.
  \item \textbf{No specialization:}
    There is no $u_0 \in \CC$ and $c \in \CC^*$ such that
    $R(u_0) = c \cdot \pisort$.
  \item \textbf{No limit:}
    There is no function $f\colon \CC \to \CC^*$ and no
    $u_0 \in \CC \cup \{\infty\}$ such that
    $\lim_{u \to u_0} R(u)/f(u) = \pisort$.
  \item \textbf{No conjugation:}
    There is no $M \in \mathrm{GL}_4(\CC)$, $u_0 \in \CC$, and
    $c \in \CC^*$ such that $M\, R(u_0)\, M^{-1} = c \cdot \pisort$.
\end{enumerate}
\end{theorem}

\begin{corollary}[Integrability Hierarchy]\label{cor:hierarchy}
The three-level hierarchy
\[
  \textit{Solvable} \;\supset\; \textit{Quasi-solvable}
  \;\supset\; \textit{Combinatorial}
\]
consists of genuinely distinct algebraic objects. The quasi-solvable level
($\pisort$ satisfying 0-Hecke relations in $H_n(0)$) is not an intermediate
degeneration of the solvable level ($R(u)$ satisfying the parametric
Yang--Baxter equation).
\end{corollary}

%% -----------------------------------------------------------------------
\section{Proof}
%% -----------------------------------------------------------------------

\subsection{Spectral data of $R(u)$}

\begin{lemma}\label{lem:block}
$R(u)$ is block-diagonal with respect to the spin decomposition
$V \otimes V = V_0 \oplus V_1 \oplus V_2$, where
$V_0 = \mathrm{span}\{\ket{00}\}$,
$V_1 = \mathrm{span}\{\ket{01}, \ket{10}\}$,
$V_2 = \mathrm{span}\{\ket{11}\}$.
\begin{itemize}
  \item On $V_0$: eigenvalue $\lambda_1 = 1+u$.
  \item On $V_2$: eigenvalue $\lambda_2 = 1$.
  \item On $V_1$: the restricted matrix is
    $\bigl(\begin{smallmatrix} u & 1 \\ 1 & 0 \end{smallmatrix}\bigr)$
    with characteristic polynomial $\lambda^2 - u\lambda - 1$, giving
    eigenvalues
    \begin{equation}\label{eq:lambdapm}
      \lambda_\pm = \frac{u \pm \sqrt{u^2 + 4}}{2}.
    \end{equation}
\end{itemize}
The four eigenvalues of $R(u)$ are therefore
$\{1+u,\; 1,\; \lambda_+,\; \lambda_-\}$.
\end{lemma}

\begin{proof}
The block structure follows from the ice rule (spin conservation).
The characteristic polynomial of the $V_1$ block is
$\det\bigl(\begin{smallmatrix} u-\lambda & 1 \\ 1 & -\lambda
\end{smallmatrix}\bigr) = \lambda^2 - u\lambda - 1$.
\end{proof}

\begin{lemma}\label{lem:det}
$\det R(u) = -(1+u)$.
\end{lemma}

\begin{proof}
$\det R(u) = (1+u) \cdot 1 \cdot
\det\bigl(\begin{smallmatrix} u & 1 \\ 1 & 0 \end{smallmatrix}\bigr)
= (1+u)(0-1) = -(1+u)$.
\end{proof}

\begin{lemma}\label{lem:ev1}
For $u \neq 0$, the eigenvalue $1$ is not a root of
$\lambda^2 - u\lambda - 1$.
\end{lemma}

\begin{proof}
Substituting $\lambda = 1$: $1 - u - 1 = -u \neq 0$.
\end{proof}

\begin{lemma}\label{lem:ev1pu}
For $u \neq 0$, the eigenvalue $1+u$ is not a root of
$\lambda^2 - u\lambda - 1$.
\end{lemma}

\begin{proof}
Substituting $\lambda = 1+u$:
$(1+u)^2 - u(1+u) - 1 = 1 + 2u + u^2 - u - u^2 - 1 = u \neq 0$.
\end{proof}

\begin{proposition}\label{prop:minpoly}
For $u \neq 0$, the minimal polynomial of $R(u)$ is
\begin{equation}\label{eq:minpoly}
  m_{R(u)}(\lambda) = (\lambda - (1+u))(\lambda^2 - u\lambda - 1)(\lambda - 1),
\end{equation}
which has degree~$4$.
\end{proposition}

\begin{proof}
The minimal polynomial is the least common multiple of the minimal
polynomials on each invariant subspace:
$(\lambda-(1+u))$ on $V_0$,
$(\lambda^2 - u\lambda - 1)$ on $V_1$ (the full characteristic polynomial,
since the $2 \times 2$ block is not scalar for $u \neq 0$),
and $(\lambda - 1)$ on $V_2$.
By Lemmas~\ref{lem:ev1} and~\ref{lem:ev1pu}, no pair shares a common
root for $u \neq 0$, so their product is their least common multiple,
giving degree $1 + 2 + 1 = 4$.
\end{proof}

%% -----------------------------------------------------------------------
\subsection{Part (1): $\pisort$ is a 0-Hecke generator}

\textbf{Idempotency.}
Direct computation shows $\pisort^2 = \pisort$. Row~3 of $\pisort$ is
identically zero, so the only nontrivial entry to check is
$(\pisort^2)_{23} = (\pisort)_{22} \cdot (\pisort)_{23} = 1 \cdot 1 = 1$.
All entries match. (Verified symbolically.)

\textbf{Braid relation.}
On $V^{\otimes 3}$ with basis $\{\ket{abc} : a,b,c \in \{0,1\}\}$,
the operators $(\pisort)_{12}$ and $(\pisort)_{23}$ sort adjacent pairs
of bits. Both compositions $\pi_{12}\pi_{23}\pi_{12}$ and
$\pi_{23}\pi_{12}\pi_{23}$ implement the full \emph{bubble sort} on
three bits, sending any triple to its nondecreasing rearrangement.
The nontrivial cases are:
\begin{center}
\begin{tabular}{@{}cccc@{}}
  \toprule
  Input & $\pi_{12}\pi_{23}\pi_{12}$ & $\pi_{23}\pi_{12}\pi_{23}$ & Output \\
  \midrule
  $\ket{100}$ &
    $\to\ket{010}\to\ket{001}\to\ket{001}$ &
    $\to\ket{100}\to\ket{010}\to\ket{001}$ &
    $\ket{001}$ \\
  $\ket{110}$ &
    $\to\ket{110}\to\ket{101}\to\ket{011}$ &
    $\to\ket{101}\to\ket{011}\to\ket{011}$ &
    $\ket{011}$ \\
  $\ket{101}$ &
    $\to\ket{011}\to\ket{011}\to\ket{011}$ &
    $\to\ket{101}\to\ket{011}\to\ket{011}$ &
    $\ket{011}$ \\
  \bottomrule
\end{tabular}
\end{center}
The remaining five inputs are already sorted and fixed by both sides.
(Verified by $8 \times 8$ matrix computation.) \qed

%% -----------------------------------------------------------------------
\subsection{Part (2): $R(u)$ is not a Hecke generator for $u \neq 0$}

A matrix satisfies a quadratic relation $(R - \alpha I)(R - \beta I) = 0$
if and only if its minimal polynomial has degree at most~2. By
Proposition~\ref{prop:minpoly}, for $u \neq 0$ the minimal polynomial of
$R(u)$ has degree~4. Therefore no quadratic relation holds. \qed

\begin{remark}
At $u = 0$, the $R$-matrix is the permutation (swap) $P$ with $P^2 = I$
and eigenvalues $\{1,1,1,-1\}$. This satisfies the quadratic
$(R-I)(R+I) = 0$, making it a generator of the symmetric group algebra
$\CC[S_n]$---the \emph{combinatorial} level of the hierarchy.
\end{remark}

%% -----------------------------------------------------------------------
\subsection{Part (3): No specialization}

Suppose $R(u_0) = c \cdot \pisort$ for some $u_0 \in \CC$, $c \in \CC^*$.
Comparing the $(1,1)$ and $(2,2)$ entries:
\[
  R(u_0)_{11} = 1 + u_0 = c, \qquad R(u_0)_{22} = u_0 = c.
\]
Subtracting: $(1+u_0) - u_0 = 1 = 0$, a contradiction.

\begin{remark}
The obstruction is structural: the diagonal entries of $R(u)$ are $1+u$
and $u$, which differ by a constant additive shift of~$1$. Those of
$\pisort$ are both~$1$. No scalar rescaling can absorb an additive
discrepancy.
\end{remark}
\qed

%% -----------------------------------------------------------------------
\subsection{Part (4): No limit}

We show that for no $u_0 \in \CC \cup \{\infty\}$ and no function $f(u)$
does $R(u)/f(u) \to \pisort$ as $u \to u_0$.

The key observation: $R(u)$ has entries depending on $u$ (namely
$R_{11} = 1+u$ and $R_{22} = u$) and constant entries ($R_{23} = R_{32} = 1$,
$R_{44} = 1$, $R_{33} = 0$). All nonzero entries of $\pisort$ equal~$1$.

\textbf{Case 1:} Finite $u_0$, $f(u) \to f_0 \neq 0$. Then
$R(u)/f(u) \to R(u_0)/f_0$, reducing to Part~(3). No solution.

\textbf{Case 2:} Finite $u_0$, $f(u) \to 0$. The constant entry
$R_{23} = 1$ gives $1/f(u) \to \infty$. Diverges.

\textbf{Case 3:} Finite $u_0$, $f(u) \to \infty$. Then
$R_{23}/f(u) \to 0$, but $(\pisort)_{23} = 1 \neq 0$. Mismatch.

\textbf{Case 4:} $u_0 = \infty$. As $u \to \infty$, the entries
$R_{11} = 1+u$ and $R_{22} = u$ grow without bound while
$R_{23} = R_{32} = 1$, $R_{44} = 1$, $R_{33} = 0$ stay bounded.
If $f(u) \sim \alpha u^\gamma$ as $u \to \infty$ for $\gamma > 0$:
\[
  \lim_{u \to \infty} \frac{R(u)}{\alpha u^\gamma} =
  \begin{cases}
    \frac{1}{\alpha}\operatorname{diag}(1,1,0,0) & \text{if } \gamma = 1, \\
    0 & \text{if } \gamma > 1, \\
    \text{diverges} & \text{if } \gamma < 1.
  \end{cases}
\]
For $\gamma = 1$, the limit has rank~2; $\pisort$ has rank~3. No match.
\qed

%% -----------------------------------------------------------------------
\subsection{Part (5): No conjugation}

This is the strongest statement. We use two conjugation invariants: the
eigenvalue multiset and the minimal polynomial.

\textbf{Step 1 (Eigenvalue constraint).}
If $M R(u_0) M^{-1} = c \cdot \pisort$, then $R(u_0)$ and
$c \cdot \pisort$ share the same eigenvalue multiset. Since $\pisort$
has eigenvalues $\{1,1,1,0\}$, the matrix $c \cdot \pisort$ has
eigenvalues $\{c,c,c,0\}$: a triple eigenvalue $c$ and a simple
eigenvalue~$0$.

\textbf{Step 2 (Where is the zero eigenvalue?)}
From the eigenvalues $\{1+u, 1, \lambda_+, \lambda_-\}$ of $R(u)$:
\begin{itemize}
  \item $1+u = 0$: $u = -1$. (Candidate.)
  \item $1 = 0$: Impossible.
  \item $\lambda_\pm = 0$: From~\eqref{eq:lambdapm},
    $\frac{u \pm \sqrt{u^2+4}}{2} = 0$ requires
    $u = \mp\sqrt{u^2+4}$, hence $u^2 = u^2+4$. Impossible.
\end{itemize}
So $u_0 = -1$ is the unique candidate.

\textbf{Step 3 (Eliminate $u_0 = -1$).}
At $u = -1$, the eigenvalues are
\begin{equation}\label{eq:evu-1}
  \Bigl\{0,\; 1,\; \frac{-1+\sqrt{5}}{2},\;
  \frac{-1-\sqrt{5}}{2}\Bigr\}.
\end{equation}
The three nonzero eigenvalues are $1$,
$\frac{-1+\sqrt{5}}{2} \approx 0.618$, and
$\frac{-1-\sqrt{5}}{2} \approx -1.618$. These are three distinct
values. For conjugacy to $c \cdot \pisort$, we need all three nonzero
eigenvalues to equal~$c$. Since they are distinct, this is impossible.

\textbf{Step 4 (Alternative: minimal polynomial argument).}
We give a second proof that works uniformly for all $u \neq 0$.
The minimal polynomial of $c \cdot \pisort$ is $\lambda(\lambda - c)$
(degree~2), since $(c\pi)^2 = c^2\pi = c(c\pi)$ and $c\pi \neq 0$,
$c\pi \neq cI$.
By Proposition~\ref{prop:minpoly}, for all $u \neq 0$ the minimal
polynomial of $R(u)$ has degree~4. The minimal polynomial is a
conjugation invariant. A matrix with minimal polynomial of degree~4
cannot be conjugate to a scalar multiple of an idempotent (minimal
polynomial degree~2).

\textbf{Step 5 (Eliminate $u_0 = 0$).}
$R(0) = P$ (the permutation matrix) with eigenvalues $\{1,1,1,-1\}$.
For conjugacy to $c \cdot \pisort$ with eigenvalues $\{c,c,c,0\}$, we
need $-1 = 0$. Impossible.

\medskip

Combining Steps 2--5: no value of $u_0 \in \CC$ admits a conjugacy
$M R(u_0) M^{-1} = c \cdot \pisort$. \qed

%% -----------------------------------------------------------------------
\section{The Conceptual Picture}
%% -----------------------------------------------------------------------

The proof reveals why the three levels of the integrability hierarchy are
algebraically incommensurable. Two independent deformation parameters
are at play:
\begin{itemize}
  \item The \textbf{spectral parameter} $u$, which parametrizes the
    $R$-matrix family $R(u)$.
  \item The \textbf{Hecke parameter} $q$, which controls the quotient
    from the braid algebra to the Hecke algebra via
    $(T_i - q)(T_i + 1) = 0$.
\end{itemize}

Setting different parameters to zero gives:

\begin{center}
\begin{tabular}{@{}llll@{}}
  \toprule
  Limit & Result & Algebra & Relation \\
  \midrule
  $u = 0$ in $R(u)$ & Permutation $P$ & $\CC[S_n]$ &
    $P^2 = I$ (involution) \\
  $q = 0$ in $H_n(q)$ & Sorting $\pisort$ & $H_n(0)$ &
    $\pi^2 = \pi$ (idempotent) \\
  \bottomrule
\end{tabular}
\end{center}

These two limits are algebraically incomparable: $P = R(0)$ is an
involution with eigenvalues $\{1,1,1,-1\}$, while $\pisort$ is an
idempotent with eigenvalues $\{1,1,1,0\}$. The sign $-1$ versus $0$
is the spectral signature of their algebraic independence.

The ``Hecke $q$'' and the ``spectral $u$'' control different aspects of
the algebraic structure. The sorting operator $\pisort$ is the $q = 0$
limit of the Hecke generator---not the $u = 0$ limit of the $R$-matrix.
These are genuinely different operations on genuinely different objects,
connected only through the braid group that is a common ancestor of both.

This connects directly to the categorical phase diagram:
\begin{itemize}
  \item \textbf{Braided monoidal} (generic $q$, generic $u$): the full
    Yang--Baxter equation holds; $R(u)$ generates a braid group action.
  \item \textbf{Coboundary monoidal} ($q = 0$): the cactus group
    replaces the braid group; $\pisort$ generates $H_n(0)$ with
    idempotent relations.
  \item \textbf{Symmetric monoidal} ($u = 0$): the permutation $P$
    generates $\CC[S_n]$ with involutive relations.
\end{itemize}
The Operator Independence Theorem says these categorical phases are
separated by algebraic walls that no continuous deformation can cross.

%% -----------------------------------------------------------------------
\section{Computational Verification}
%% -----------------------------------------------------------------------

All matrix computations were independently verified using SymPy
(symbolic) and NumPy (numerical):

\begin{itemize}
  \item \textbf{Part (1):} Idempotency $\pisort^2 = \pisort$ and braid
    relation verified by direct $4 \times 4$ and $8 \times 8$ matrix
    multiplication (\texttt{hecke\_integrability\_test.py}).
  \item \textbf{Part (2):} Eigenvalues of $R(u)$ computed symbolically;
    quadratic relation $(R-\alpha I)(R-\beta I) = 0$ tested at
    $u = 1, 2, \frac{1}{2}, 3$ with no annihilator found
    (\texttt{pi\_sort\_independence.py}).
  \item \textbf{Part (3):} Entrywise system $R(u_0) = c \cdot \pisort$
    solved symbolically---empty solution set.
  \item \textbf{Part (4):} Limits computed symbolically for
    $u \to 0, \pm 1, \infty$ with normalizations
    $f(u) = 1, u, 1+u, au+b$.
  \item \textbf{Part (5):} Jordan forms and eigenvalue multisets compared
    at all candidate points $u = 0, -1$.
\end{itemize}

\noindent
Scripts: \texttt{hecke\_integrability\_test.py},
\texttt{pi\_sort\_independence.py}.

\end{document}
